\begin{enumerate}

\item Momentos respecto a $A$: $R_B \cdot 5 = 6 \cdot 5 \cdot 2{,}5 + 12 + 10 \cdot 7$; el momento horario aplicado entra con su valor, sin brazo, y del mismo lado que las cargas, que también giran en sentido horario respecto a $A$.

\item Fuerzas verticales: $R_A + R_B = 6 \cdot 5 + 10$.

\item En el voladizo el cortante es constante e igual a la carga $P$; el flector varía linealmente hasta anularse en $C$.

\item En $D$ el flector salta 12 kN m hacia los positivos, porque un momento horario sobre la parte izquierda comprime la fibra superior a su derecha.

\item El cortante del tramo $AB$ se anula en $x = R_A/q$; ese punto queda a la izquierda de $D$, de modo que el flector máximo positivo está justo a la derecha del salto.

\end{enumerate}

\noindent\textbf{Operaciones y resultados intermedios.}

\begin{enumerate}

\item Reacción en B.

\[\fitmath{y_{1}=\frac{6\cdot 5\cdot 2.5+12+10\cdot 7}{5}\approx 31.4\,\mathrm{kN}}\]

\item Reacción en A.

\[\fitmath{y_{2}=6\cdot 5+10-31.4\approx 8.6\,\mathrm{kN}}\]

\item Flector a la derecha de D.

\[\fitmath{y_{3}=8.6\cdot 2.5-\frac{6\cdot {2.5}^{2}}{2}+12\approx 14.75\,\mathrm{kN m}}\]

\item Flector sobre el apoyo B.

\[\fitmath{y_{4}=-10\cdot 2\approx -20\,\mathrm{kN m}}\]

\end{enumerate}
